% Part: incompleteness
% Chapter: representability-in-q
% Section: minimization-representable

\documentclass[../../../include/open-logic-section]{subfiles}

\begin{document}

\olfileid{inc}{req}{min}
\olsection{Panggolekan Tanpa Wates Reguler Representabel ing $\Th{Q}$}

Ayo nimbang panggolekan tanpa wates. Upamana $g(x, z)$ reguler lan
representabel ing $\Th{Q}$, upamane dening !!{formula}~$!A_g(x, z, y)$.
Tetepna $f$ ditegesi dening $f(z) = \umin{x}{[g(x, z) = 0]}$. Kita
arep nemokake !!a{formula}~$!A_f(z, y)$ kang merepresentasikake~$f$.
Biji~$f(z)$ iku bilangan~$x$ kang (a) nyukupi $g(x, z) = 0$ lan (b)
paling cilik ing antarane bilangan kang nyukupi iki, yaiku kanggo
saben $w < x$, $g(w, z) \neq 0$. Mula pilihan kang alamiah yaiku:
\[
!A_f(z,y) \ident !A_g(y, z, \Obj 0) \land \lforall[w][(w < y \lif \lnot
  !A_g(w, z, \Obj 0))].
\]
Ing kasus umum, mesthi wae kita kudu ngganti $z$ karo
$z_0$, \dots, $z_k$.

Bukti iki maneh mbutuhake sawetara lema ngenani perkara kang
bisa dibuktekake dening $\Th{Q}$.

\begin{lem}
\ollabel{lem:succ} Kanggo saben !!{constant}~$a$ lan saben bilangan
asli~$n$,
\[
\Th{Q} \Proves \eq[(a' + \num n)][(a + \num n)'].
\]
\end{lem}

\begin{proof}
Buktine, kaya lumrahe, nganggo induksi marang $n$. Ing kasus dhasar,
$n = 0$, kita kudu nuduhake yen $\Th{Q}$ mbuktekake
$\eq[(a' + \Obj 0)][(a + \Obj 0)']$. Nanging kita nduweni:
\begin{align}
  \Th{Q} & \Proves \eq[(a' + \Obj 0)][a'] \quad \text{miturut aksioma $Q_4$}
  \ollabel{step1}\\
  \Th{Q} & \Proves  \eq[(a + \Obj 0)][a] \quad \text{miturut aksioma $Q_4$}
  \ollabel{step2} \\
  \Th{Q} & \Proves \eq[(a + \Obj 0)'][a'] \quad \text{miturut \olref{step2}}
  \ollabel{step3} \\
  \Th{Q} & \Proves \eq[(a' + \Obj 0)][(a + \Obj 0)'] \quad
  \text{miturut \olref{step1} lan \olref{step3}}\notag
\end{align}
Ing langkah induksi, kita bisa nganggep wis nuduhake yen $\Th{Q}
\Proves \eq[(a' + \num n)][(a + \num n)']$. Amarga
$\num{n+1}$ iku $\num{n}'$, kita kudu nuduhake yen $\Th{Q}$
mbuktekake $\eq[(a' + \num{n}')][(a + \num{n}')']$. Kita nduweni:
\begin{align}
  \Th{Q} & \Proves \eq[(a' + \num n')][(a' + \num n)'] \quad
  \text{miturut aksioma $!Q_5$} \ollabel{step5}\\
  \Th{Q} & \Proves \eq[(a' + \num n')][(a + \num n')'] \quad
  \text{hipotesis induksi} \ollabel{step6}\\
  \Th{Q} & \Proves \eq[(a' + \num n)'][(a + \num n')'] \quad
  \text{miturut \olref{step5} lan \olref{step6}.} \notag
\end{align}
\end{proof}

Perlu dielingake maneh yen iki luwih ringkih tinimbang nyatakake yen
$\Th{Q}$ mbuktekake $\lforall[x][\lforall[y][(x' + y) = (x + y)']]$.
Sanajan !!{sentence} iki bener ing~$\Struct{N}$, $\Th{Q}$ ora
mbuktekake.

\begin{lem}
\ollabel{lem:less-zero}
$\Th{Q} \Proves \lforall[x][\lnot x < \Obj 0]$.
\end{lem}

\begin{proof}
Kita menehi bukti kanthi informal (yaiku mung menehi pituduh
carane mbangun !!{derivation} formal).

Kita kudu mbuktekake $\lnot a < \Obj 0$ kanggo~$a$ sembarang. Miturut
definisi $<$, kita kudu mbuktekake $\lnot
\lexists[y][\eq[(y'+a)][\Obj 0]]$ ing $\Th{Q}$. Kita bakal nganggep
$\lexists[y][\eq[(y'+a)][\Obj 0]]$ lan mbuktekake kontradiksi. Upamana
$\eq[(b'+a)][\Obj 0]$. Kanthi $!Q_3$, kita nduweni
$\eq[a][\Obj 0] \lor \lexists[y][\eq[a][y']]$. Kita bedakake
kasuse.

Kasus 1: $\eq[a][\Obj 0]$ lumaku. Saka $\eq[(b'+a)][\Obj 0]$, kita entuk
$\eq[(b' + \Obj 0)][\Obj 0]$. Miturut aksioma~$!Q_4$ saka $\Th{Q}$,
kita nduweni $\eq[(b' + \Obj 0)][b']$, mula $\eq[b'][\Obj 0]$. Nanging
miturut aksioma~$!Q_2$ kita uga nduweni $\eq/[b'][\Obj 0]$,
kontradiksi.

Kasus 2: Kanggo sawijining $c$, $\eq[a][c']$. Nanging banjur kita
nduweni $\eq[(b' + c')][\Obj 0]$. Miturut aksioma~$!Q_5$, kita nduweni
$\eq[(b' + c)'][\Obj 0]$, kang maneh mbantah aksioma~$Q_2$.
\end{proof}

\begin{lem}
\ollabel{lem:less-nsucc}
Kanggo saben bilangan asli~$n$,
  \[
  \Th{Q} \Proves
  \lforall[x][(x < \num {n+1} \lif (\eq[x][\Obj 0] \lor \dots \lor
    \eq[x][\num n]))].
  \]
\end{lem}

\begin{proof}
Kita nggunakake induksi marang~$n$. Ayo nimbang kasus dhasar, nalika
$n = 0$. Ing kasus iki, kita kudu nuduhake $a < \num 1 \lif
\eq[a][\Obj 0]$ kanggo~$a$ sembarang. Upamana $a < \num 1$. Mula
miturut aksioma kang netepake $<$, kita nduweni
$\lexists[y][\eq[(y'+a)][\Obj 0']]$ (amarga $\num 1 \ident
\Obj 0'$).

Upamana $b$ nduweni sipat iku, yaiku $\eq[(b'+a)][\Obj 0']$. Kita
kudu nuduhake $\eq[a][\Obj 0]$. Miturut aksioma~$!Q_3$, kita nduweni
$\eq[a][\Obj 0]$ utawa ana $c$ saengga $\eq[a][c']$. Ing kasus
kapisan, ora ana maneh kang kudu dituduhake. Mula upamana
$\eq[a][c']$. Banjur kita nduweni $\eq[(b' + c')][\Obj 0']$.
Miturut aksioma~$!Q_5$ saka $\Th{Q}$, kita nduweni
$\eq[(b'+c)'][\Obj 0']$. Miturut aksioma~$!Q_1$, kita nduweni
$\eq[(b' + c)][\Obj 0]$. Nanging miturut aksioma~$!Q_8$ iki ateges
$c < \Obj 0$, kang mbantah \olref{lem:less-zero}.

Saiki langkah induksi. Kita mbuktekake kasus kanggo $n+1$ kanthi
nganggep kasus kanggo~$n$. Upamana $a < \num {n+2}$. Kanthi $!Q_3$
maneh, kita bisa mbedakake rong kasus: $\eq[a][\Obj 0]$ lan ana $b$
saengga $\eq[a][b']$. Ing kasus kapisan, $\eq[a][\Obj 0] \lor \dots
\lor \eq[a][\num{n+1}]$ langsung lumaku. Ing kasus kapindho, kita
nduweni $b' < \num {n+2}$, yaiku $b' < \num{n+1}'$. Miturut aksioma
~$!Q_8$, kanggo sawijining $c$, $\eq[(c'+b')][\num{n+1}']$. Miturut
aksioma $!Q_5$, $\eq[(c'+b)'][\num{n+1}']$. Miturut aksioma~$!Q_1$,
$\eq[(c'+b)][\num{n+1}]$, mula $b < \num{n+1}$ miturut aksioma~$!Q_8$.
Miturut hipotesis induksi, $\eq[b][\Obj 0] \lor \dots \lor
\eq[b][\num{n}]$. Saka kene, kita entuk $\eq[b'][\Obj 0'] \lor
\dots \lor \eq[b'][\num{n}']$ kanthi logika, mula
$\eq[a][\num{1}] \lor \dots \lor \eq[a][\num{n+1}]$ amarga
$\eq[a][b']$.
\end{proof}

\begin{lem}
  \ollabel{lem:trichotomy} For every natural number~$m$,
  \[
  \Th{Q} \Proves
  \lforall[y][((y < \num{m} \lor \num{m} < y) \lor \eq[y][\num{m}])].
  \]
\end{lem}

\begin{proof}
Kanthi induksi marang~$m$. Luwih dhisik, nimbang kasus $m=0$.
$\Th{Q} \Proves
\lforall[y][(\eq[y][\Obj 0] \lor \lexists[z][\eq[y][z']])]$ miturut~$!Q_3$.
Tetepna $a$ sembarang. Mula $\eq[a][\Obj 0]$ utawa ana~$b$ saengga
$\eq[a][b']$. Ing kasus kapisan, kita uga nduweni $(a < \Obj 0 \lor
\Obj 0 < a) \lor \eq[a][\Obj 0]$. Nanging yen $\eq[a][b']$, mula
$\eq[(b' + \Obj 0)][(a + \Obj 0)]$ miturut logika~$\eq$. Miturut
$!Q_4$, $\eq[(a + \Obj 0)][a]$, mula kita nduweni
$\eq[(b' + \Obj 0)][a]$, lan banjur
$\lexists[z][\eq[(z' + \Obj 0)][a]]$. Miturut definisi $<$ ing
$!Q_8$, $\Obj 0 < a$. Yen $\Obj 0 < a$, mula uga
$(\Obj 0 < a \lor a < \Obj 0) \lor \eq[a][\Obj 0]$.

Saiki upamana kita nduweni
\begin{align*}
  \Th{Q} & \Proves \lforall[y][((y < \num{m} \lor \num{m} < y) \lor
    \eq[y][\num{m}])]
  \intertext{lan kita arep nuduhake}
  \Th{Q} & \Proves \lforall[y][((y < \num{m+1} \lor \num{m+1} < y) \lor
    \eq[y][\num{m+1}])]
\end{align*}
Tetepna $a$ sembarang. Miturut $!Q_3$, $\eq[a][\Obj 0]$ utawa kanggo
sawijining~$b$, $\eq[a][b']$. Ing kasus kapisan, kita nduweni
$\eq[\num{m}' + a][\num{m+1}]$ miturut $!Q_4$, mula $a < \num{m+1}$
miturut $!Q_8$.

Saiki nimbang kasus kapindho, $\eq[a][b']$. Miturut hipotesis
induksi, $(b < \num{m} \lor \num{m} < b) \lor
    \eq[b][\num{m}]$.

Disjungsi kapisan $b < \num{m}$ padha tegese (miturut $!Q_8$) karo
$\lexists[z][\eq[(z' + b)][\num{m}]]$. Upamana $c$ nduweni sipat iki.
Yen $\eq[(c' + b)][\num{m}]$, mula uga
$\eq[(c' + b)'][\num{m}']$. Miturut $!Q_5$,
$\eq[(c' + b)'][(c' + b')]$. Mula $\eq[(c' +
b')][\num{m}']$. Kita entuk
$\lexists[u][\eq[(u' + b')][\num{m+1}]]$ kanthi ngenalake kuantor
eksistensial marang~$c'$ lan ngelingi yen $\num{m}' \ident
\num{m+1}$. Dadi yen $b < \num{m}$, mula $b' < \num{m+1}$ lan
$a < \num{m+1}$.

Saiki upamana $\num{m} < b$, yaiku $\lexists[z][\eq[(z' +
\num{m})][b]]$. Upamana $c$ iku~$z$ kang dimaksud, yaiku
$\eq[(c' + \num{m})][b]$. Kanthi logika,
$\eq[(c' + \num{m})'][b']$. Miturut $!Q_5$,
$\eq[(c' + \num{m}')][b']$. Amarga $\eq[a][b']$ lan
$\num{m}' \ident \num{m+1}$, kita entuk
$\eq[(c' + \num{m+1})][a]$. Miturut $!Q_8$, $\num{m+1} < a$.

Pungkasan, upamana $\eq[b][\num{m}]$. Mula kanthi logika,
$\eq[b'][\num{m}']$, dadi $\eq[a][\num{m+1}]$.

Mula saka saben disjungsi ing kasus kanggo~$m$ lan~$b$, kita bisa
entuk disjungsi kang cocog kanggo~$m+1$ lan~$a$.
\end{proof}
  
\begin{prop}
\ollabel{prop:rep-minimization}
Yen $!A_g(x, z, y)$ merepresentasikake $g(x, z)$ ing~$\Th{Q}$, mula
\[
!A_f(z,y) \ident !A_g(y, z, \Obj 0) \land \lforall[w][(w < y \lif \lnot
  !A_g(w, z, \Obj 0))]
\]
merepresentasikake $f(z) = \umin{x}{[g(x, z) = 0]}$.
\end{prop}

\begin{proof}
Luwih dhisik, kita nuduhake yen $f(n) = m$, mula
$\Th{Q} \Proves !A_f(\num n, \num m)$, yaiku,
\begin{align}
  \Th{Q} & \Proves !A_g(\num{m}, \num{n}, \Obj 0) \land \lforall[w][(w <
    \num{m} \lif \lnot !A_g(w, \num{n}, \Obj 0))]. \notag
  \intertext{Amarga $!A_g(x, z, y)$ merepresentasikake $g(x, z)$ lan
    $g(m, n) = 0$ yen $f(n) = m$, kita nduweni}
\Th{Q} & \Proves !A_g(\num{m}, \num{n}, \Obj 0). \notag
\intertext{Yen $f(n) = m$, mula kanggo saben $k < m$, $g(k, n) \neq 0$.
  Dadi}
\Th{Q} & \Proves \lnot !A_g(\num{k}, \num{n}, \Obj 0). \notag
\intertext{Kita entuk}
\Th{Q} & \Proves \lforall[w][(w < \num{m} \lif \lnot
  !A_g(w, \num{n}, \Obj 0))]. \ollabel{rep-less}
\end{align}
miturut \olref{lem:less-zero} yen $m = 0$ lan miturut
\olref{lem:less-nsucc} yen ora.

Saiki ayo nuduhake yen $f(n) = m$, mula $\Th{Q} \Proves
\lforall[y][(!A_f(\num{n}, y) \lif \eq[y][\num{m}])]$. Kita maneh
nggambarake argumentasi kanthi informal, lan ninggalake formalisasine
kanggo pamaca.

Upamana $!A_f(\num{n}, b)$. Saka iki kita entuk (a)
$!A_g(b, \num{n}, \Obj 0)$ lan (b)
$\lforall[w][(w < b \lif \lnot !A_g(w, \num{n}, \Obj
  0))]$. Miturut \olref{lem:trichotomy},
$(b < \num{m} \lor \num{m} < b) \lor \eq[b][\num{m}]$. Kita bakal
nuduhake yen $b < \num{m}$ lan $\num{m} < b$ loro-lorone njalari
kontradiksi.

Yen $\num{m} < b$, mula $\lnot !A_g(\num{m}, \num{n}, \Obj 0)$
saka~(b). Nanging $m = f(n)$, mula $g(m, n) = 0$, lan
$\Th{Q} \Proves !A_g(\num{m}, \num{n}, \Obj 0)$ amarga $!A_g$
merepresentasikake~$g$. Dadi ana kontradiksi.

Saiki upamana $b < \num{m}$. Amarga
$\Th{Q} \Proves \lforall[w][(w <
  \num{m} \lif \lnot !A_g(w, \num{n}, \Obj 0))]$ miturut
\olref{rep-less}, kita entuk $\lnot !A_g(b, \num{n}, \Obj 0)$.
Iki maneh mbantah~(a).
\end{proof}

\end{document}
